% Final Report — due Dec 8, 5 pages, 11 pt.
% Required sections: introduction, related works, methodology, experimental
% results, conclusions.
% Build: make final
\documentclass[11pt]{article}
\usepackage{deqreport}
\reporttitle{Final Report}

\begin{document}
\makeheader

\section{Introduction}

% Open with the problem and the finding, not with background. One paragraph on
% what a DEQ is, one on what this project did, one on what was found.
\textsc{TODO}

\section{Related Work}

% Reuse and expand the proposal's related work. Group by theme, not by paper.
\paragraph{Implicit layers.}
\textsc{TODO}~\cite{amos2017optnet,chen2018neural,elghaoui2021implicit}

\paragraph{Deep equilibrium models.}
\textsc{TODO}~\cite{bai2019deq,bai2020mdeq,winston2020monotone}

\paragraph{Cheaper and more stable training.}
\textsc{TODO}~\cite{bai2021stabilizing,geng2021training,fung2022jfb}

\section{Methodology}

\subsection{Equilibrium formulation}

For input $x$ the hidden state is the fixed point
$\zstar = \fth(\zstar, x)$, and the gradient follows from implicit
differentiation:
\begin{equation}
  \dd{\ell}{\theta}
  = \dd{\ell}{\zstar}\left(I - \dd{\fth}{\zstar}\right)^{-1}\dd{\fth}{\theta}.
  \label{eq:implicit}
\end{equation}

\subsection{Forward solver}

% Include the Anderson pseudocode — it is your implementation, so show it.
\begin{algorithm}[h]
\caption{Anderson acceleration for $\zstar = \fth(\zstar,x)$}
\label{alg:anderson}
\begin{algorithmic}[1]
  \State \textsc{TODO: fill from deq/solvers.py}
\end{algorithmic}
\end{algorithm}

\subsection{Backward pass}

% Derive the hand-computed Jacobian-vector products. This is the core
% contribution for the "implement backpropagation yourself" track.
\textsc{TODO}

\subsection{Gradient approximations}

\textsc{TODO: exact, Neumann order $k$, Jacobian-free.}

\section{Experimental Results}

\subsection{Setup}

% Datasets, model sizes, optimizer, epochs, hardware. Enough to reproduce.
\textsc{TODO}

\subsection{Gradient correctness}

\textsc{TODO: the finite-difference check from Report I.}

\subsection{Accuracy, memory and time}

\begin{table}[h]
  \centering
  \caption{\textsc{TODO}}
  \begin{tabularx}{\textwidth}{@{}Xrrrr@{}}
    \toprule
    Method & Test acc.\ (\%) & Peak mem.\ (MB) & Time/epoch (s) & Grad.\ cos.\ sim. \\
    \midrule
    \textsc{TODO} & & & & \\
    \bottomrule
  \end{tabularx}
  \label{tab:main}
\end{table}

\subsection{Ablations}

% Solver tolerance, Anderson memory, Neumann order.
\textsc{TODO}

\section{Conclusions}

% What was learned, what the trade-off is, and one honest limitation.
\textsc{TODO}

\printrefs
\end{document}
